\section{Problem Setting}
\label{sec:problem-setting}

Let $D$ denote a fitted dataframe with rows $x_i$ and columns partitioned into
continuous variables $N$ and categorical variables $C$. The task considered in
this paper is not general-purpose private data release, and it is not
unconditional density estimation over all possible mixed-type tables. Instead,
we consider the constructive problem of fitting an inspectable sampler to $D$
and producing a new dataframe $\tilde D$ whose rows are useful for tests,
demonstrations, notebooks, and exploratory analyses. The output must therefore
remain in the original column schema, while the intermediate representation may
be numerical.

The type policy is part of the problem definition. A column is treated as
continuous only when it has a numeric dtype and more than two distinct
non-missing values. Boolean columns, string columns, and two-valued numeric
columns are treated as categorical, because a stored binary code is usually a
label rather than a continuous coordinate. This convention is conservative: it
prevents binary labels from entering Euclidean geometry as if the difference
between $0$ and $1$ were a calibrated numeric distance.

A fitted sampler is evaluated under three constraints. First, it should avoid
trivial row reuse or synthetic rows that are closer to training rows than real
rows are to one another under ordinary nearest-neighbour diagnostics. Second,
it should preserve enough marginal, dependence, and task-relevant information
to be useful in low-stakes example-data workflows. Third, its failure modes
should be visible through fitted components and diagnostics. These constraints
define a practical-use target rather than a guarantee: if formal privacy,
regulatory release, or hard domain constraints are required, they must be added
outside the sampler.
